\chapter{Methodology and Proposed Framework}
\label{chap:methodology}

% -----------------------------------------------------------------------
% TO WRITE: Describe the overall design of your comparison framework and
% the technical choices behind each method. Explain *why* you chose each
% approach, what data you used, and how you prepared it.
% Suggested structure below — remove these comments as you write.
% -----------------------------------------------------------------------

\section{Research Design and Framework Overview}
% Describe the overall structure of your comparison:
% - What are the four methods being compared and why those four?
% - How does the experiment flow (historical simulation first, then live)?
% - Provide a diagram/figure of the framework pipeline if possible.

\section{Dataset Description}
% Describe the ProvenExpert campaign dataset:
% - Time range, number of campaigns, number of recipients
% - What features are available (device, geo, engagement history, etc.)
% - Any preprocessing, anonymization, or filtering steps applied

\section{Campaign Design Parameters}
\label{sec:campaign_parameters}
% List the email elements being optimized:
% - Subject line variants
% - CTA button variants
% - Layout / content block variants
% - Image placement / color scheme variants
% Explain how variants were generated or selected.

\section{Baseline: Classical A/B Testing}
% Describe your A/B testing implementation:
% - Traffic split method (random assignment)
% - Statistical test used (e.g. chi-squared, z-test)
% - Stopping rule (fixed sample size or significance threshold)

\section{Bayesian Bandit with Thompson Sampling}
% Describe the Thompson Sampling implementation:
% - Prior distribution chosen (e.g. Beta-Bernoulli for CTR)
% - How posterior is updated with each new observation
% - How variant selection decisions are made per send

\section{Contextual Bandit with User Segmentation}
% Describe the contextual bandit implementation:
% - Which user features are used as context
% - Algorithm used (e.g. LinUCB or contextual Thompson Sampling)
% - How context vectors are constructed per recipient

\section{Reinforcement Learning for Drip Campaigns}
% Describe the RL formulation:
% - State space (user engagement history, campaign step)
% - Action space (next email variant to send)
% - Reward function (e.g. click, conversion, or composite)
% - Algorithm used (e.g. Q-learning, DQN, policy gradient)
% - Training procedure on historical data

\section{Evaluation Metrics and Criteria}
% Define formally how you will compare the four methods:
% - Primary metrics: CTR, open rate, conversion rate
% - Secondary metrics: convergence speed, regret (for bandits)
% - Segment-level breakdown (device, region, engagement tier)
